\subsection{Re-Scaling and Re-Parameterizing to Integer Counts}\label{subsec:alg_rescaling}

%\textbf{Instructional Time:} 2 hours

The DG algorithm\index{Doob-Gillespie algorithm} simulates discrete jumps in state by $\pm 1$ at a time, so we scale\index{rescaling} the ODE system such that the state variable discretely counts individuals (such as molecules or organisms). This requires a change of variables.\index{change of variables}


\begin{example}\label{ex:radioactive_decay_units}
Take a fictional isotope with measured half-life $\tau = 5.00 \times 10^{3}\,\mathrm{s}$ and initial sample mass $M_0 = 2.5\,\mathrm{lb}$. Your boss demands a simulation running in years and kilograms. The per-second decay rate is
\[
  \mu = \frac{\ln 2}{\tau}
      = \frac{\ln 2}{5.00 \times 10^{3}\,\mathrm{s}}
      \approx 1.386 \times 10^{-4}\,\mathrm{s}^{-1}.
\]
The relevant conversion factors are
\[
  1\,\mathrm{lb} = 0.4536\,\mathrm{kg},
  \qquad
  1\,\mathrm{yr} = 3.156 \times 10^{7}\,\mathrm{s}.
\]
Whereas $M_0$ is measured in pounds and $\tau$ is measured in seconds, let $\widetilde{M}$ denote mass in kilograms and $T$ time in years.
Substituting $M = \widetilde{M}/0.4536$ and $t = T \cdot 3.156 \times 10^{7}$
into $\frac{dM}{dt} = -\mu M$ gives
\[
  \frac{1}{0.4536} \cdot \frac{1}{3.156 \times 10^{7}}
  \frac{d\widetilde{M}}{dT}
  =
  -\mu \cdot \frac{\widetilde{M}}{0.4536}.
\]
The factor $1/0.4536$ appears on both sides and cancels, leaving
\[
  \frac{d\widetilde{M}}{dT}
  = -\underbrace{\bigl(\mu \cdot 3.156 \times 10^{7}\bigr)}_{\displaystyle\widetilde{\mu}}\,
  \widetilde{M}
  \approx -4374\,\mathrm{yr}^{-1} \cdot \widetilde{M}.
\]
The initial condition converts as
$\widetilde{M}_0 = 2.5 \times 0.4536 \approx 1.134\,\mathrm{kg}$.
Because the ODE is linear in mass, and with the same unit on both sides,
the mass conversion factor cancels. Only the time unit requires a numerical
change to $\mu$.
\end{example}

\begin{example}\label{ex:saline_tank}
\textbf{Saline Tank Model.} A nanomachinery manufacturer has made a motivating observation for mathematical modeling: a sensitive manufacturing process requires a saline solution with precise quantities of salt to avoid failure. The manufacturer articulates the problem in the following way.
\begin{itemize}
\item If too many or too few NaCl molecules (really, Na-Cl ion pairs in aqueous solution) enter the mixing tank, the industrial process fails, so we must model the absolute number of molecules.
\item A tank with $V$ units of volume, two inlet pipes, and one outlet pipe holds the well-mixed saline solution.
\item One inlet pipe brings in distilled water at rate $a$ units of volume per unit of time.
\item Another inlet pipe brings in a saline solution at $b$ units of volume per unit of time with concentration $c$ units of mass per unit of volume.
\item The outlet pipe drains $a+b$ units of volume of the well-mixed solution per unit of time.
\end{itemize}

We formulate an ODE governing the salt concentration in the tank as $\frac{dX}{dt} = \frac{bc}{V} - \frac{a+b}{V}X$. The saline inlet adds $bc$ units of mass per unit time to a tank of volume $V$, so the concentration increases at rate $bc/V$. The outlet removes the well-mixed solution at total volumetric rate $a+b$, diluting the concentration at rate $(a+b)X/V$.

The number of molecules $M$ times mass per molecule $m$ is the total mass, and the total mass per volume $V$ is the concentration. Let $X = mM/V$, where $M$ is the number of NaCl molecules in the tank, $m$ is the mass of each molecule in the same mass units as $X$ and $c$, and $V$ is the volume. Taking the derivative of both sides of $X = mM/V$ gives us the following.
\begin{align}
\frac{m}{V}M =& X \notag \\
\frac{d}{dt} \frac{m}{V}M =& \frac{d}{dt} X \notag \\
\frac{m}{V} \frac{dM}{dt} =& \frac{dX}{dt} \notag \\
\frac{m}{V} \frac{dM}{dt} =& \frac{bc}{V} - \frac{a+b}{V}X \notag \\
\frac{m}{V} \frac{dM}{dt} =& \frac{bc}{V} - \frac{a+b}{V} \frac{mM}{V} \notag \\
m \frac{dM}{dt} =& bc - m\frac{a+b}{V}M \notag \\
\frac{dM}{dt} =& \frac{bc}{m} - \frac{a+b}{V}M \label{eqn:saline_tank}
\end{align} Now the evolution of $M$ takes place on the scale of individual molecules. This process works regardless of the underlying parameters $a$, $b$, $c$, $m$, and $V$.

\end{example}

Sometimes a change of variables requires changing time as in Example~\ref{ex:radioactive_decay_units} or Homework Problem~\ref{ex:hw_brusselator_transform}
Sometimes, different system parameters are given in different unit systems (e.g., a volume measured in liters but a flow rate in gallons per second).
Example~\ref{ex:saline_tank_with_repar} illustrates unit-conversion with mixed units.
When mixed units are on the table, we can get real weird with it.



\begin{example}\label{ex:saline_tank_with_repar}
\textbf{Saline Tank Model with Mixed Units.}
Mixed units occur frequently during meta-analyses comparing data from multiple sources, measured in units distinct to each source. Convert all variables to compatible unit systems before applying the methods in this chapter.

Consider the tank in the previous problem, except the parameters $V$, $a$, $b$, $c$, and $m$ appear in inconsistent units: say $V$ is in liters, $a$ and $b$ are both in cubic light-years per year, $c$ is in kilograms per gallon, and $m$ is in pounds. Because the units are inconsistent, the concentration of salt in the tank no longer obeys the differential equation $\frac{dX}{dt} = \frac{bc}{V} - \frac{a+b}{V}X$ and the mass does not obey $\frac{dM}{dt} = \frac{bc}{m} - \frac{a+b}{V}M$. The term $bc/V$ adds a quantity with units $\frac{(\text{light-year})^3}{\text{year}} \cdot \frac{\text{kg}}{\text{gal}} \cdot \frac{1}{\text{liter}}$, which one cannot add to $\frac{a+b}{V}X$, measured in different units. Do not perform arithmetic on quantities with mismatched units.


Reformulate the model using liters, hours, and grams, with new parameters $\widetilde{V}$, $\widetilde{a}$, $\widetilde{b}$, $\widetilde{c}$, and $\widetilde{m}$.
\begin{itemize}
\item The parameter $V$ is already in liters, so set $\widetilde{V} = V$.
\item The parameters $a$ and $b$ are both in cubic light-years per year,
$\frac{(\textup{light year})^3}{year}$,
which is a measure of volume over time. We need these parameters in liters per hour,
$\frac{\textup{liter}}{\textup{hour}}$
. Readers should verify for themselves that there are $9.667\times 10^{46} \frac{\textup{liter}}{\textup{hour}}$ in
$1 \frac{(\textup{light year})^3}{\textup{year}}$.
So, $a \frac{(\textup{light year})^3}{\textup{year}} =  9.667 a \times 10^{46} \frac{\textup{liter}}{\textup{hour}}$. Set $\widetilde{a} = 9.667a \times 10^{46} \frac{\textup{liter}}{\textup{hour}}$ and, similarly, $\widetilde{b} = 9.667b\times 10^{46} \frac{\textup{liter}}{\textup{hour}}$.
\item The parameter $c$ is in kilograms per gallon $\frac{\textup{kg}}{\textup{gal}}$. There are $264.2$ grams per liter $\frac{\textup{gram}}{\textup{liter}}$ in $1 \frac{\textup{kilogram}}{\textup{gallon}}$. So, set $\widetilde{c} = 264.2c \frac{\textup{gram}}{\textup{liter}}$.
\item The parameter $m$ is in pounds $\textup{lbs}$ per molecule. There are $453.6\textup{ grams}$ per $1\textup{ lbs}$. So, set $\widetilde{m} = 453.6m \frac{\textup{grams}}{\textup{molecule}}$.
\end{itemize} Now the parameters $\widetilde{V}$, $\widetilde{a}$, $\widetilde{b}$, $\widetilde{c}$, and $\widetilde{m}$ can replace $V, a, b, c, m$ as in Example~\ref{ex:saline_tank}, giving $\frac{dM}{dt} = \frac{\widetilde{b}\widetilde{c}}{\widetilde{m}} - \frac{\widetilde{a}+\widetilde{b}}{\widetilde{V}}M$. The equation has the same algebraic form as before. Only the numerical values of the parameters have changed for consistency in units.



\end{example}

\textup{{\bf Questions for Reflection:}}
\begin{enumerate}
\item In Example~\ref{ex:radioactive_decay_units}, the factor $1/0.4536$ appeared on both sides of the rescaled ODE and canceled. Explain in general why a multiplicative rescaling of the state variable does not change the qualitative dynamics of a linear ODE, but does change the numerical value of the rate parameter.

\item The saline tank model introduces the mass per molecule $m$ to convert mass to molecule count. As $m \to 0$ with total mass held fixed, the molecule count grows without bound. Describe what happens to the DG simulation in this limit and why the continuous ODE remains a valid description.

\item A colleague proposes running a DG simulation with parameters measured in inconsistent units, arguing that the results will look reasonable enough. Construct a specific counterexample using the saline tank model to show that unit inconsistency produces a numerically incorrect simulation, and identify the symptom a researcher would observe.
\end{enumerate}
